\section{Interpretation Plan}
\label{sec:discussion}

The study is designed to distinguish several possible outcomes rather than
confirm a preferred legitimacy narrative.

\begin{itemize}
\item A benchmark advantage over enacted maps could reflect trust in
reproducibility, dislike of legislatures, visual differences, treatment
wording, or demand effects.
\item Equivalence with commission maps would not show that algorithms and
commissions are democratically interchangeable.
\item A transparency effect would be normatively useful only if respondents
also understand that the benchmark does not guarantee fairness, VRA
compliance, community preservation, or partisan symmetry.
\item Partisan heterogeneity could be larger among political elites or when
respondents evaluate their own districts.
\item Null or negative results would be informative evidence about adoption
barriers and must be reported.
\end{itemize}

The design must also test adversarial messages that accurately describe
parameter discretion, geographic sorting, legal modifications, and known
evidence gaps. A treatment that works only when these limitations are hidden
would not support a durable legitimacy claim.

Results from two states and online panels will not establish national public
acceptance. External replication, state-specific studies, and real-stakes
deliberative settings would remain necessary.
